%%%%%%%%%%%%%%%%%%%%%%%%%%%%%%%%%%%%%%%%%%%%%%%%%%%%%%%%%
%
%    Name            ::  AGH University of Krakow theme for LaTeX Beamer presentations
%    Author          ::  Stanisław Polak (polak[aT]icsr[DoT]agh[DoT]edu[DoT]pl)
%    Created         ::  2017-03-30
%    Modified        ::  2024-03-18
%    Version         ::  2.0
%    Email           ::  polak[aT]icsr[DoT]agh[DoT]edu[DoT]pl
%    Website         ::  http://www.icsr.agh.edu.pl/~polak/
%    Youtube         ::  https://www.youtube.com/user/spolak69
%    Github          ::  https://github.com/polaksta
%    License         ::  This work may be distributed and/or modified under the
%                            conditions of the LaTeX Project Public License, either version 1.3
%                            of this license or (at your option) any later version.
%                        The latest version of this license is in
%                          http://www.latex-project.org/lppl.txt
%                          and version 1.3 or later is part of all distributions of LaTeX
%                          version 2005/12/01 or later.
%
%    Description     ::    This presentation is a demonstration of the 'AGH' beamer theme.
%                          It is  based on the standard beamer theme 'CambridgeUS' 
%                              created by Madhusudan Singh
%                    
%
%    Theme           ::    dark
%    Options         ::    dark = darknesslevel
%                    ::    parttitle = author | date | title | institute
%                    ::    outertheme = default | infolines | miniframes | smoothbars | sidebar | split | shadow | tree | smoothtree
%                    ::    outerthemeoptions = <options>  
%                    ::    margins = <value>  
%                    ::    nosidebar                   
%%%%%%%%%%%%%%%%%%%%%%%%%%%%%%%%%%%%%%%%%%%%%%%%%%
\documentclass[
    english,          % Jeżeli chcesz otrzymać prezentację w języku polskim, to wyraz 'english' zastąp wyrazem 'polish' 
    aspectratio=1610, % Proporcje ekranu to 16:10
    % handout,        % Odkomentuj, jeśli chcesz przełączyć z trybu 'beamer' na 'handout'
]
{beamer}
\usepackage{babel}
\makeatletter
\@ifclasswith{beamer}{polish}{
    \usepackage{polski}
}
\makeatother
\usepackage[utf8]{inputenc}
\usepackage{listings} 
\usepackage{minted}   
 %%%%%%%%% Ustawienia dla trybu 'beamer' %%%%%%%%%%%%%%
\mode<beamer>{ 
    \definecolor{links}{HTML}{2A1B81}
    \hypersetup{
        pdfpagemode=FullScreen, % Włącz tryb pełnoekranowy
        colorlinks,
        linkcolor=,
        urlcolor=links
    }
   \usetheme[parttitle=date]{AGH}                                                             % Zamiast daty (w stopce slajdu) wyświetlaj tytuł bieżącej części 
    %\usetheme[nosidebar]{AGH}                                                                % Nie pokazuj paska bocznego na nietytułowych slajdach
    %\usetheme[nosidebar, margins=1em]{AGH}                                                   % Nie pokazuj paska bocznego na nietytułowych slajdach, ustaw obydwa marginesy (lewy/prawy) na wartość 1em
    %\usetheme[dark]{AGH}                                                                     % Użyj ciemnego tła
    %\usetheme[dark, parttitle=author]{AGH}                                                   % Użyj ciemnego tła, a w stopce, zamiast imienia i nazwiska autora, wyświetlaj tytuł części
    %\usetheme[outertheme=miniframes, outerthemeoptions={footline=authorinstitutetitle}]{AGH} % Użyj motywu zewnętrznego 'miniframes' z opcją 'footline=authorinstitutetitle'
}
 %%%%%%%%% Ustawienia dla trybu 'handout' %%%%%%%%%%%%%%
\mode<handout>{
    \hypersetup{pdfpagemode=UseNone}
    \usepackage{pgfpages}
    \pgfpagesuselayout{4 on 1}[a4paper,border shrink=5mm,landscape] % Pokaż 4 slajdy na 1 stronie
    \pgfpageslogicalpageoptions{1}{border code=\pgfusepath{stroke}}
    \pgfpageslogicalpageoptions{2}{border code=\pgfusepath{stroke}}
    \pgfpageslogicalpageoptions{3}{border code=\pgfusepath{stroke}}
    \pgfpageslogicalpageoptions{4}{border code=\pgfusepath{stroke}}
    \usetheme{boxes}
    \addheadbox{structure}{\quad\insertpart\hfill\insertsection\hfill\insertsubsection\qquad}          % Treść nagłówka
    \addfootbox{structure}{\quad\insertshortauthor\hfill\insertframenumber\hfill\insertsubtitle\qquad} % Treść stopki
}
 %%%%%%%%%%%%%%%%%%%%%%%%%%%%%%%%%%%%%%%%%%%%%%%%%%%%%%
\AtBeginPart{
    \frame{\partpage} % Na początku każdej części będzie, automatycznie, wstawiany slajd z tytułem (części)
}
 %%%%%%%%%%%%%%%%%%%%%%%%%%%%%%%%%%%%%%%%%%%%%%%%%%%%%%
\date{} % Nie wyświetlaj (bieżącej) daty na slajdzie tytułowym
 %%%%%%%%%%%%%%%%%%%%%%%%%%%%%%%%%%%%%%%%%%%%%%%%%%%%%%
\titlegraphic{
    \includegraphics[width=1.5cm]{img/csg-logo.png} % Grafika widoczna na slajdzie tytułowym
}                               
 %%%%%%%%%%%%%%%%%%%%%%%%%%%%%%%%%%%%%%%%%%%%%%%%%%%%%%
\logo{
    \includegraphics[width=2cm]{img/logo_podstawowe.png} % Grafika widoczna ma każdym ze slajdów
}                       
 %%%%%%%%%%%%%%%%%%%%%%%%%%%%%%%%%%%%%%%%%%%%%%%%%%%%%%
\author[Stanisław Polak]{Stanisław Polak 
 % \texorpdfstring{\newline{\scriptsize \href{mailto:polak@agh.edu.pl}{polak@agh.edu.pl}}}{}    % Odkomentuj, jeżeli, poniżej imienia i nazwiska, ma być widoczny adres e-mail
}
 %%%%%%%%%%%%%%%%%%%%%%%%%%%%%%%%%%%%%%%%%%%%%%%%%%%%%%
\iflanguage{polish}{
    \title{Temat AGH dla Beamer{\dywiz}a}
    \subtitle{Przykładowa prezentacja}
    \institute[AGH]{
            Wydział Informatyki\newline
            ul. Kawiory 21\newline
            30-055 Kraków\newline
            \url{http://www.icsr.agh.edu.pl/~polak/}
   }
}{
    \title{Beamer theme AGH}
    \subtitle{Sample presentation}
    \institute[AGH]{
            Faculty of Computer Science\newline
            Kawiory 21 Street\newline
            30-055 Kraków\newline
            Poland\newline
            \url{http://www.icsr.agh.edu.pl/~polak/}
   }
}
 %%%%%%%%%%% Konfiguracja pakietu 'listings' %%%%%%%%%%%%%%%%%%%%%%%%%%
 % Źródło: https://en.wikibooks.org/wiki/LaTeX/Source_Code_Listings#Using_the_listings_package
\lstset{
    showstringspaces=false,                           % Czy podkreślać spacje wewnątrz napisów?   
    numbers=left,                                     % Gdzie umieścić numery linii?   
    stepnumber=1,                                     % Krok pomiędzy dwoma numerami linii. Jeśli 1, każda linia będzie numerowana
    numbersep=5pt,                                    % Jak daleko numery linii są odsunięte od kodu
    showspaces=false,                                 % Czy wyświetlać spacje w ciągu znaków przy użyciu podkreśleń?
    showtabs=false,                                   % Czy wyświetlać tabulatory, w napisach, przy użyciu podkreśleń?
    frame=single,                                     % Dodaj pojedynczą ramkę wokół kodu
    tabsize=2,                                        % Ustaw domyślny rozmiar tabulatora na dwie spacje
    captionpos=t,                                     % Ustaw pozycję podpisu — podpisy na górze
    breaklines=true,                                  % Włącz automatyczne łamanie linii
    breakatwhitespace=false,                          % Czy automatyczne łamanie powinno odbywać się tylko w przypadku białych znaków?
    keywordstyle=\ttfamily\color{blue},               % Styl słów kluczowych
    identifierstyle=\ttfamily\color{violet}\bfseries, % Styl identyfikatora
    commentstyle=\color{brown},                       % Styl komentarzy
    stringstyle=\ttfamily,                            % Styl napisów
    escapeinside={\%*}{*)},                           % Tekst znajdujący się pomiędzy znacznikiem '%*' a znacznikiem '*)' to komendy LaTeX — kompilator będzie je wykonywał
    % Pakiet listings nie obsługuje polskich znaków UTF-8, ale można to „naprawić” — polskie znaki zostaną zastąpione poleceniami LaTeX — patrz poniżej    
    inputencoding=utf8,
    literate={ą}{{\k{a}}}1
             {Ą}{{\k{A}}}1
             {ę}{{\k{e}}}1
             {Ę}{{\k{E}}}1
             {ó}{{\'o}}1
             {Ó}{{\'O}}1
             {ś}{{\'s}}1
             {Ś}{{\'S}}1
             {ł}{{\l{}}}1
             {Ł}{{\L{}}}1
             {ż}{{\.z}}1
             {Ż}{{\.Z}}1
             {ź}{{\'z}}1
             {Ź}{{\'Z}}1
             {ć}{{\'c}}1
             {Ć}{{\'C}}1
             {ń}{{\'n}}1
             {Ń}{{\'N}}1
}
 %%%%%%%%%%% Konfiguracja pakietu 'minted' %%%%%%%%%%%%
\setminted[C++]{
    frame=single, % Pojedyncze obramowanie wokół kodu
     linenos      % Numeruj linie kodu
} 
 %%%%%%%%%%%%%%%%%%%%%%%%%%%%%%%%%%%%
 %  % Konfiguracja rysunków TikZ
 %%%%%%%%%%%%%%%%%%%%%%%%%%%%%%%%%%%%
\usetikzlibrary{shapes.geometric,arrows.meta} 
\tikzstyle{every picture}+=[remember picture]
\tikzstyle{na} = [baseline=-.5ex]
 %%%%%%%%%%%%%%%%%%%%%%%%%%%%%%%%%%%%%%%%%%%%%%%%%%%%%%
\begin{document}
 %%%%%%%%%%%%%%%%%%%%%%%%%%%%%%%%%%%%%%%%%%%%%%%%%%%%%%
 %%%%%%%%%%%%%%%%%% Slajd tytułowy %%%%%%%%%%%%%%%%%%%%
 %%%%%%%%%%%%%%%%%%%%%%%%%%%%%%%%%%%%%%%%%%%%%%%%%%%%%%
\maketitle
 %%%%%%%%%%%%%%%%%%%%%%%%%%%%%%%%%%%%%%%%%%%%%%%%%%%%%%
 %%%%%%%%%%%%%%%%%% Slajd 1 %%%%%%%%%%%%%%%%%%%%%%%%%%%
 %%%%%%%%%%%%%%%%%%%%%%%%%%%%%%%%%%%%%%%%%%%%%%%%%%%%%%
\ifdefined\textleftmargin 
  \begin{frame}[fragile]{\iflanguage{polish}{Informacje}{Information}}
    \begin{center}
      \iflanguage{polish}{
        $\Longleftarrow$ Aktualna wartość \hfill Aktualna wartość $\Longrightarrow$\\
        $\Longleftarrow$ rozmiaru lewego marginesu \hfill rozmiaru prawego marginesu $\Longrightarrow$\\
        $\Longleftarrow$ to \the\textleftmargin \hfill to \the\textrightmargin $\Longrightarrow$
      }{
        $\Longleftarrow$ The current value of \hfill The current value of $\Longrightarrow$\\
        $\Longleftarrow$ the  left  margin size \hfill the right margin size $\Longrightarrow$\\
        $\Longleftarrow$ is \the\textleftmargin \hfill is \the\textrightmargin $\Longrightarrow$
      }
    \end{center}
    \iflanguage{polish}{
      Możesz je zmieniać za pomocą parametru 'margins' \pauza
    }{
      You can change them with the 'margins' parameter ---
    }
    \verb+\usetheme[margins=...]{AGH}+
  \end{frame}
\fi
 %%%%%%%%%%%%%%%%%%%%%%%%%%%%%%%%%%%%%%%%%%%%%%%%%%%%%%   
 %%%%%%%%%%%%%%%%%% Część „Przykłady” %%%%%%%%%%%%%%%%%
 %%%%%%%%%%%%%%%%%%%%%%%%%%%%%%%%%%%%%%%%%%%%%%%%%%%%%%   
\part{\iflanguage{polish}{Przykłady}{Examples}}
 %%%%%%%%%%%%%%%%%%%%%%%%%%%%%%%%%%%%%%%%%%%%%%%%%%%%%%
 %%%%%%%%%%%%%%%%%% Slajd 2 %%%%%%%%%%%%%%%%%%%%%%%%%%%
 %%%%%%%%%%%%%%%%%%%%%%%%%%%%%%%%%%%%%%%%%%%%%%%%%%%%%%
\begin{frame}{\iflanguage{polish}{Plan prezentacji}{Outline}}
    % Plan prezentacji jest tworzony na podstawie treści komend '\section' oraz '\subsection' zawartych w bieżącej części
    \tableofcontents[pausesections] % Utwórz „przyrostowy” (elementy odkrywane po kolei) plan prezentacji
\end{frame}
 %%%%%%%%%%%%%%%%%%%%%%%%%%%%%%%%%%%%%%%%%%%%%%%%%%%%%%   
 %%%%%%% Sekcja „Elementy podstawowe” %%%%%%%%%%%%%%%%%
 %%%%%%%%%%%%%%%%%%%%%%%%%%%%%%%%%%%%%%%%%%%%%%%%%%%%%%  
\section{\iflanguage{polish}{Elementy podstawowe}{Basic elements}}
 %%%%%%%%%%%%%%%%%%%%%%%%%%%%%%%%%%%%%%%%%%%%%%%%%%%%%%
 %%%%%%%%%%%%%%%%%% Slajd 3 %%%%%%%%%%%%%%%%%%%%%%%%%%%
 %%%%%%%%%%%%%%%%%%%%%%%%%%%%%%%%%%%%%%%%%%%%%%%%%%%%%%
\begin{frame}{\iflanguage{polish}{Wyszczególnienie}{Itemize}}
  \begin{columns}
    \column{0.5\textwidth} % Pierwsza kolumna
    \begin{itemize}
      \item \iflanguage{polish}{Element 1}{Item 1}
      \item \iflanguage{polish}{Element 2}{Item 2}
      \item \iflanguage{polish}{Element 3}{Item 3}
    \end{itemize}
    \column{0.5\textwidth} % Druga kolumna
    \pause
    \structure{\iflanguage{polish}{Odkrywanie po kolei}{Uncovering one by one}}
    \begin{itemize}[<+->]
      \item \iflanguage{polish}{Element 1}{Item 1}
      \item \iflanguage{polish}{Element 2}{Item 2}
      \item \iflanguage{polish}{Element 3}{Item 3}
    \end{itemize}
    \onslide
  \end{columns}
\end{frame}
 %%%%%%%%%%%%%%%%%%%%%%%%%%%%%%%%%%%%%%%%%%%%%%%%%%%%%%
 %%%%%%%%%%%%%%%%%% Slajd 4 %%%%%%%%%%%%%%%%%%%%%%%%%%%
 %%%%%%%%%%%%%%%%%%%%%%%%%%%%%%%%%%%%%%%%%%%%%%%%%%%%%%
\begin{frame}{\iflanguage{polish}{Wyliczenie}{Enumerate}}
  \begin{columns}
    \column{0.5\textwidth}
    \begin{enumerate}
      \item \iflanguage{polish}{Element 1}{Item 1}
      \item \iflanguage{polish}{Element 2}{Item 2}
      \item \iflanguage{polish}{Element 3}{Item 3}
    \end{enumerate}
    \column{0.5\textwidth}
    \pause
    \structure{\iflanguage{polish}{Odkrywanie elementów po kolei z jednoczesnym wyróżnianiem}{Uncovering elements in turn with simultaneous highlighting}}
    \begin{enumerate}[<+-|alert@+>]
      \item \iflanguage{polish}{Element 1}{Item 1}
      \item \iflanguage{polish}{Element 2}{Item 2}
      \item \iflanguage{polish}{Element 3}{Item 3}
    \end{enumerate}
    \onslide
  \end{columns}
\end{frame}
 %%%%%%%%%%%%%%%%%%%%%%%%%%%%%%%%%%%%%%%%%%%%%%%%%%%%%%   
 %%%%%%% Sekcja „Matematyka” %%%%%%%%%%%%%%%%%%%%%%%%%%
 %%%%%%%%%%%%%%%%%%%%%%%%%%%%%%%%%%%%%%%%%%%%%%%%%%%%%%  
\section{\iflanguage{polish}{Matematyka}{Mathematics}}
 %%%%%%%%%%%%%%%%%%%%%%%%%%%%%%%%%%%%%%%%%%%%%%%%%%%%%%
 %%%%%%%%%%%%%%%%%% Slajd 5 %%%%%%%%%%%%%%%%%%%%%%%%%%%
 %%%%%%%%%%%%%%%%%%%%%%%%%%%%%%%%%%%%%%%%%%%%%%%%%%%%%%
\begin{frame}{\iflanguage{polish}{Podstawowe bloki}{Basic blocks}}
  % Przykłady z podręcznika użytkownika klasy Beamer
  \iflanguage{polish}{
    \begin{block}{Definicja}
      \alert{Zbiór} składa się z elementów.
    \end{block}
    \begin{exampleblock}{Przykład}
      Zbiór $\{1,2,3,5\}$ zawiera cztery elementy.
    \end{exampleblock}
    \begin{alertblock}{Błędne Twierdzenie}
      $1=2$.
    \end{alertblock}
  }{
    \begin{block}{Definition}
      A \alert{set} consists of elements.
    \end{block}
    \begin{exampleblock}{Example}
      The set $\{1,2,3,5\}$ has four elements.
    \end{exampleblock}
    \begin{alertblock}{Wrong Theorem}
      $1=2$.
    \end{alertblock}
  }
\end{frame}
 %%%%%%%%%%%%%%%%%%%%%%%%%%%%%%%%%%%%%%%%%%%%%%%%%%%%%%
 %%%%%%%%%%%%%%%%%% Slajd 6 %%%%%%%%%%%%%%%%%%%%%%%%%%%
 %%%%%%%%%%%%%%%%%%%%%%%%%%%%%%%%%%%%%%%%%%%%%%%%%%%%%%
\begin{frame}{\iflanguage{polish}{Otoczenia matematyczne}{Math environments}}
  \begin{columns}
    \column{0.45\textwidth} 
    \structure{\iflanguage{polish}{Twierdzenia}{Theorems}}
    \begin{theorem}[\iflanguage{polish}{Pitagorasa}{Pythagorean}]
      $a^{2}+  b^{2}=  c^{2}$
    \end{theorem}
    \column{0.45\textwidth}
    \structure{\iflanguage{polish}{Dowody}{Proofs}}
    \begin{proof}
      \ldots
    \end{proof}
  \end{columns}
  \vfill
  \ldots
  \vfill
  \begin{definition}
    \ldots
  \end{definition}
\end{frame}
 %%%%%%%%%%%%%%%%%%%%%%%%%%%%%%%%%%%%%%%%%%%%%%%%%%%%%%
 %%%%%%%%%%%%%%%%%% Slajd 7 %%%%%%%%%%%%%%%%%%%%%%%%%%%
 %%%%%%%%%%%%%%%%%%%%%%%%%%%%%%%%%%%%%%%%%%%%%%%%%%%%%%
\begin{frame}{\iflanguage{polish}{Dynamiczny wzór matematyczny}{Dynamic mathematical formula}}
  \[
    \binom{n}{k} = \pause \frac{n!}{k!(n-k)!}
  \]
\end{frame}
 %%%%%%%%%%%%%%%%%%%%%%%%%%%%%%%%%%%%%%%%%%%%%%%%%%%%%%
 %%%%%%%%%%%%%%%%%% Slajd 8 %%%%%%%%%%%%%%%%%%%%%%%%%%%
 %%%%%%%%%%%%%%%%%%%%%%%%%%%%%%%%%%%%%%%%%%%%%%%%%%%%%%
\begin{frame}{\iflanguage{polish}{Rysowanie po slajdzie}{Drawing on the slide}}
\iflanguage{polish}{Ułamek składa się z:}{Every fraction consists of:}
\begin{itemize}
\item<2-> \iflanguage{polish}{Licznika}{Numerator} \tikz[na] \node[coordinate] (n1) {};
\end{itemize}

\begin{displaymath}
a =  \frac{
	\tikz[baseline]{
	    \node[fill=blue!20,anchor=base] (t1){$x+y$};
        }
      }
      {
        \tikz[baseline]{
            \node[fill=red!20, ellipse,anchor=base] (t2){$y-z$};
        }
      }
\end{displaymath}

\begin{itemize}
\item<3> \iflanguage{polish}{Mianownika}{Denominator} \tikz[na]\node [coordinate] (n2) {};
\end{itemize}
\begin{tikzpicture}[overlay]
        \path[->]<2-> (n1) edge [bend left] (t1);
        \path[->]<3-> (n2) edge [bend right] (t2);
\end{tikzpicture}
\end{frame}
 %%%%%%%%%%%%%%%%%%%%%%%%%%%%%%%%%%%%%%%%%%%%%%%%%%%%%%   
 %%%%%%% Sekcja „Informatyka” %%%%%%%%%%%%%%%%%%%%%%%%%
 %%%%%%%%%%%%%%%%%%%%%%%%%%%%%%%%%%%%%%%%%%%%%%%%%%%%%%  
\section{\iflanguage{polish}{Informatyka}{Computer Science}}
 %%%%%%%%%%%%%%%%%%%%%%%%%%%%%%%%%%%%%%%%%%%%%%%%%%%%%%   
 %%%%%%% Podsekcja „Wstawianie kodów źródłowych” %%%%%%%
 %%%%%%%%%%%%%%%%%%%%%%%%%%%%%%%%%%%%%%%%%%%%%%%%%%%%%%  
\subsection*{\iflanguage{polish}{Wstawianie kodów źródłowych}{Inserting source codes}}
 %%%%%%%%%%%%%%%%%%%%%%%%%%%%%%%%%%%%%%%%%%%%%%%%%%%%%%
 %%%%%%%%%%%%%%%%%% Slajd 9 %%%%%%%%%%%%%%%%%%%%%%%%%%%
 %%%%%%%%%%%%%%%%%%%%%%%%%%%%%%%%%%%%%%%%%%%%%%%%%%%%%%
\begin{frame}[fragile]{\iflanguage{polish}{Użycie otoczenia 'listings'}{Using the 'listings' environment}}
  \begin{lstlisting}[language=C++]
/* The first  program in C++ */  %*\pause*)
#include <iostream>  %*\pause*)
using namespace std; %*\pause*)
void main()
{       %*\pause*)
  cout %*\pause*) << "Hello World!"%*\pause*) << endl; %*\onslide<4->*)
} %*\onslide*)
\end{lstlisting}
\end{frame}
 %%%%%%%%%%%%%%%%%%%%%%%%%%%%%%%%%%%%%%%%%%%%%%%%%%%%%%
 %%%%%%%%%%%%%%%%%% Slajd 10 %%%%%%%%%%%%%%%%%%%%%%%%%%
 %%%%%%%%%%%%%%%%%%%%%%%%%%%%%%%%%%%%%%%%%%%%%%%%%%%%%%
\begin{frame}[fragile]{\iflanguage{polish}{Użycie otoczenia 'minted'}{Using the 'minted' environment}}
  \begin{minted}[beameroverlays, escapeinside=||]{C++}
/* The first  program in C++ */  |\pause|
#include <iostream>  |\pause|
using namespace std; |\pause|
void main()
{       |\pause|
  cout |\pause| << "Hello World!"|\pause| << endl; |\onslide<4->|
} |\onslide|
    \end{minted}
\end{frame}
 %%%%%%%%%%%%%%%%%%%%%%%%%%%%%%%%%%%%%%%%%%%%%%%%%%%%%%
 %%%%%%%%%%%%%%%%%% Dodatek %%%%%%%%%%%%%%%%%%%%%%%%%%%
 %%%%%%%%%%%%%%%%%%%%%%%%%%%%%%%%%%%%%%%%%%%%%%%%%%%%%%
\appendix
 %%%%%%%%%%%%%%%%%%%%%%%%%%%%%%%%%%%%%%%%%%%%%%%%%%%%%%
 %%%%%%%%%%%%%%%%%% Slajd 11 %%%%%%%%%%%%%%%%%%%%%%%%%%
 %%%%%%%%%%%%%%%%%%%%%%%%%%%%%%%%%%%%%%%%%%%%%%%%%%%%%%
\begin{frame}[fragile]{\iflanguage{polish}{Informacje}{Information}}
  \iflanguage{polish}{Aktualna wersja szablonu jest dostępna pod adresem \url{https://github.com/polaksta/LaTeX/tree/master/beamerthemeAGH}}{The current version of the template is available at \url{https://github.com/polaksta/LaTeX/tree/master/beamerthemeAGH}}
\end{frame}
 %%%%%%%%%%%%%%%%%%%%%%%%%%%%%%%%%%%%%%%%%%%%%%%%%%%%%%
 %%%%%%%%%%%%%%%%%% Slajd 12 %%%%%%%%%%%%%%%%%%%%%%%%%%
 %%%%%%%%%%%%%%%%%%%%%%%%%%%%%%%%%%%%%%%%%%%%%%%%%%%%%%
 % 'allowframebreaks' — jeżeli zawartość nie mieści się na jednym slajdzie, to automatycznie utwórz kolejne
\begin{frame}[allowframebreaks]{\iflanguage{polish}{Bibliografia}{Bibliography}}
  \begin{thebibliography}{9}
    %%%%%%%%%%%%%%%%%%%
    % Punktor 'online'
    %%%%%%%%%%%%%%%%%%%
    \setbeamertemplate{bibliography item}[online]  
    \bibitem{wikibook}{Wikibooks \newblock \LaTeX/Source Code Listings \newblock \url{https://en.wikibooks.org/wiki/LaTeX/Source_Code_Listings}}
    \bibitem{beamer}{Till Tantau, Joseph Wright, Vedran Miletić \newblock The beamer class \newblock \url{http://mirror.ctan.org/macros/latex/contrib/beamer/doc/beameruserguide.pdf}}
    %%%%%%%%%%%%%%%%%%%
    % Punktor 'książka'
    %%%%%%%%%%%%%%%%%%%
    \setbeamertemplate{bibliography item}[book]   
    \bibitem{lamport}{Leslie Lamport \newblock LATEX: a document preparation system : user's guide and reference manual \newblock Addison-Wesley Pub. Co., 1994 }
    %%%%%%%%%%%%%%%%%%%
    % Punktor 'artykuł'
    %%%%%%%%%%%%%%%%%%%
    \setbeamertemplate{bibliography item}[article] 
    \iflanguage{polish}{
      \bibitem{article1}{Autor \newblock Tytuł artykułu \newblock Edytor, rok \newblock Uwagi}
      %%%%%%%%%%%%%%%%%%%
       % Punktor 'trójkąt'
       %%%%%%%%%%%%%%%%%%%
      \setbeamertemplate{bibliography item}[triangle]
      \bibitem{article2}{Autor \newblock Tytuł artykułu \newblock Edytor, rok \newblock Uwagi}
      %%%%%%%%%%%%%%%%%%%
      % Punktory tekstowe
      %%%%%%%%%%%%%%%%%%%
      \setbeamertemplate{bibliography item}[text]
      \bibitem{article3}{Autor \newblock Tytuł artykułu \newblock Edytor, rok \newblock Uwagi}
      \bibitem[Polak98]{article4}{Autor \newblock Tytuł artykułu \newblock Edytor, rok \newblock Uwagi}
    }{
      \bibitem{article1}{Author \newblock Title of the article\newblock Editor, year \newblock Notes}
      %%%%%%%%%%%%%%%%%%%
      % Punktor 'trójkąt'
      %%%%%%%%%%%%%%%%%%%
      \setbeamertemplate{bibliography item}[triangle]
      \bibitem{article2}{Author \newblock Title of the article\newblock Editor, year \newblock Notes}
      %%%%%%%%%%%%%%%%%%%
      % Punktory tekstowe
      %%%%%%%%%%%%%%%%%%%
      \setbeamertemplate{bibliography item}[text]
      \bibitem{article3}{Author \newblock Title of the article\newblock Editor, year \newblock Notes}
      \bibitem[Polak98]{article4}{Author \newblock Title of the article\newblock Editor, year \newblock Notes}
    }
  \end{thebibliography}
\end{frame}
\end{document}